% sections/supplement_body.tex -- supplementary material of the MMHM paper:
% setup details, noise, detailed validation, complete results and the
% secondary analyses (experiment identifiers E0--E8 are those of
% experiments/configs/). Measured numbers only via \Num macros and generated
% tables. References to the main paper carry the prefix main- (xr).

\section{Setup details}\label{sec:s-setup}

This supplement accompanies the main paper. Sections, tables and figures
numbered with an S are in this document; all other theorems, sections, tables
and figures are those of the main paper. The experiments are organised in the
configuration files by the identifiers E0--E8 of \cref{tab:sweeps}; this
supplement uses them to point to the data.

\subsection{Implementation and platform}\label{sec:exp-platform}

\paragraph{Engine.}
All methods except \method{gudhi} run in one C++17 program and share its
complex store, which maps the sorted vertex tuple of a simplex to its facets and
cofaces and validates every atomic edit. Order keys are dense 32-bit integers; a
(re)build assigns $0,\dots,N-1$ in the chosen order, and every later insertion
takes the next unused key (Section~\ref{main-sec:drb}). A chain is a strictly increasing
array of keys, so its low is the last entry and a $\Z$ addition is a linear
merge; a working column that receives many short chains switches to a max-heap
with lazy cancellation. Deletions therefore pay the sorted-array cost of
Proposition~\ref{main-prop:cost}, and insertions its ordered-set cost up to a logarithmic
factor. The footprint index keeps the exact count $|\occ(\sigma)|$ and a lazily
filtered list of slots. In Case~A the engine chooses the cycle $V_z$ with the
fewest cells (ties: the smallest slot identifier). Only the methods that may
compute a Morse order after $K_0$ keep an index of the cells by vertex tuple;
among the DRB methods that is \method{mmhm}. The collapse of
Definition~\ref{main-def:collapse} uses a first-in-first-out queue initialised in (dimension,
vertex tuple) order, so the matching depends only on the complex. Debug builds
check (I1)--(I3), the pivot table and the footprint index after every operation.

\paragraph{Platform.}
All runs use one cloud virtual machine with an Intel Xeon processor at
2.10\,GHz (4~cores, one thread per core; 48\,KiB L1 data and 2\,MiB L2 cache per
core, 260\,MiB shared L3 cache) and 15\,GiB of memory, running Ubuntu 24.04.5
LTS with Linux~6.18. The engine is compiled by g++~13.3.0 (CMake 3.28.3) with
\texttt{-O2 -DNDEBUG} and no machine-specific tuning. The GUDHI baseline and
oracle run under Python~3.13.16 with GUDHI~3.13.0 and numpy~2.5.3. Every method
runs single-threaded.

\subsection{Further notes on the methods}\label{sec:exp-methods}

\paragraph{Static baselines.}
\method{full-id}, \method{full-sorted} and \method{full-morse} recompute all
ranks after every update by the static reduction with clearing
(Section~\ref{main-sec:drb}), without storing~$V$: \method{full-id} in insertion order,
\method{full-sorted} in the canonical order, \method{full-morse} in the Morse
order of a collapse matching computed anew after every update. By
Corollaries~\ref{main-cor:critical-work} and~\ref{main-cor:morse-reduction} the column additions of
\method{full-morse} all occur in critical columns and, after the gradient-flow
steps, are those of the reduction of the Morse complex, but it pays for a
collapse after every update. The insertion order of
\method{full-id} is the order in which the cells now present were inserted, so
it depends on the history of edits: refilled cells receive late keys
(\cref{sec:exp-timing} reports the consequence). The canonical order of
\method{full-sorted} depends only on the current complex, so it cannot drift
with the history; it is our history-independent reference. All three share the
store and the reduction code with the dynamic methods, so differences reflect
mostly the algorithms; the remaining implementation differences are this
history dependence and the index of cells by vertex tuple that
\method{full-sorted} and \method{mmhm} maintain (E4).
\method{gudhi} replays the same edits on a GUDHI simplex
tree~\cite{maria2014gudhi} through its Python interface and, after every update,
calls \texttt{compute\_persistence} with the arguments
\texttt{homology\_coeff\_field=2} and \texttt{persistence\_dim\_max=True}, and then
\texttt{betti\_numbers}; without the second argument GUDHI omits the top
dimension. GUDHI computes persistence pairs, a richer output, and is not designed
for this task; it serves as a reference point for a widely used library and as
an oracle.

\paragraph{Dynamic methods.}
\method{mmhm} is the method of Sections~\ref{main-sec:drb} and~\ref{main-sec:morse}: a dynamic reduced basis built
in the Morse order and rebuilt in a fresh Morse order under the rebuild policy of
Section~\ref{main-sec:drb}, with $m=\infty$ and $\alpha=1$ unless stated otherwise (E7
varies both). Three ablations separate its ingredients. \method{drb-id} uses
insertion keys and never rebuilds. \method{mmhm-norb} uses the Morse order of
$K_0$ and never rebuilds. \method{drb-id-rb} rebuilds under the same policy and
parameters as \method{mmhm}, but its rebuild reruns the static reduction with
clearing in the current insertion order (the keys in use, renumbered) instead
of computing a new matching. So \method{mmhm} against \method{drb-id-rb}
isolates the order used at a rebuild, and \method{mmhm} against
\method{mmhm-norb} isolates rebuilding. \method{drb-rand} builds the DRB in a
uniformly random order of the cells of each dimension, drawn from the seed, and
never rebuilds; against \method{drb-id} and \method{mmhm-norb} it tests whether
a structured order matters at all (E8).

\paragraph{Surface counters.}
\method{euler} maintains $\chi$ and the numbers of edges with one, two or
more cofaces in $O(1)$ per atomic edit and returns the Betti numbers of
Remark~\ref{main-rem:surfaces}. At the end of every update it checks that the complex has a
triangle, no tetrahedron, and no edge outside one or two triangles, and aborts
otherwise; connectivity and the absence of pinched vertices are assumed, not
checked.


\subsection{Workload generators}\label{sec:exp-workloads}


Five generators produce the update sequences. Each is
deterministic given its parameters and seed, emits only valid atomic edits
(checked again by an independent validator), lists deletions top-down and
insertions bottom-up, and records its parameters, the sizes of all states and
the number of edits per update. For W1, W2, W3 and W5 the Betti numbers of every
state are known in closed form and are written next to the sequence. Distances
are graph distances in the base mesh.

\paragraph{W1: punctured surfaces.}
The base surface is the geodesic icosphere of frequency $f=2^L$, with $10f^2+2$
vertices, $30f^2$ edges and $20f^2$ triangles, or the $n\times m$ periodic grid
torus with every square split along one diagonal, with $nm$ vertices, $3nm$
edges and $2nm$ triangles. For a vertex $v$, the patch $P(v,k)$ is the set of
triangles incident to a vertex at distance at most $k-1$ from~$v$; we use $k=2$.
When $v$ and its neighbours have degree six, $P(v,2)$ has $24$ triangles and
opening it takes $61$ atomic edits; no vertex of our surfaces has a larger
degree, so no update has more edits. A patch around a
random vertex is used only if its closure is a disk and its vertices are at
distance more than one from the vertices of every existing hole. Opening a hole
deletes the triangles of the patch, then its edges, then the vertices left
without cofaces; closing a hole is the exact reverse of its opening. With $h$
holes, the generator opens a hole if $h=0$; otherwise it opens one with
probability $0.75$, $0.5$ or $0.25$ when $h$ is below, at or above the target
$3$, and closes one otherwise or when no valid patch is found. The Betti numbers
are $(1,\max(h-1,0),[h=0],0)$ for the sphere and $(1,2+\max(h-1,0),[h=0],0)$ for
the torus, where $[h=0]$ is $1$ if $h=0$ and $0$ otherwise.
W1 is the setting of Remark~\ref{main-rem:surfaces}, which E5 examines. It also contains
the phenomenon of Proposition~\ref{main-prop:nonlocal}: when no hole is open, the surface is a
connected closed surface, so the deletion that opens the first hole costs every
DRB method $\Omega(n_2)$ operations by Proposition~\ref{main-prop:nonlocal}(a), and the last
triangle inserted when the last hole closes costs $\Omega(n_2)$ by
Proposition~\ref{main-prop:nonlocal}(b).

\paragraph{W2: refinement and coarsening in a 3-ball.}
The base complex is the Freudenthal (Kuhn) triangulation of a grid of
$N\times N\times N$ cubes, each split into six tetrahedra; it has $(N+1)^3$
vertices and about $26N^3$ simplices. A refinement applies $s$ stellar $1\to4$
subdivisions to $s$ tetrahedra collected breadth-first over face adjacency from
a uniformly random tetrahedron; tetrahedra created by earlier subdivisions may
be subdivided again. A coarsening undoes $s$ leaf subdivisions, those whose four
tetrahedra are all still present: a uniformly random one and the $s-1$ nearest
to it, measured by the Euclidean distance between the inserted vertices. A
subdivision and its inverse each consist of $16$ atomic edits, so $s_t=16s$.
Because coarsening can only undo subdivisions, $K_0$ already contains $8s$ live
subdivisions, and with $x$ live subdivisions the generator refines with
probability $\min\bigl(0.95,\max(0.05,\,0.5-0.5(x-8s)/(8s))\bigr)$, and always
when fewer than $s$ leaf subdivisions exist. Since $|K|=|K_{\mathrm{base}}|+14x$,
the size performs a random walk around its initial value. The Betti
numbers are $(1,0,0,0)$ after every update, but not within an update: by
Theorem~\ref{main-thm:exact}(d) every atomic edit changes one Betti number by one.

\paragraph{W3: cavities in a 3-ball.}
The base complex is that of W2. A centre $v$ is drawn uniformly among the
interior vertices, and the ball $B(v,k)$ is the set of tetrahedra incident to a
vertex at distance at most $k-1$ from~$v$; we use $k=2$. Carving deletes the
tetrahedra of $B$ and then, top-down, every face left without cofaces. A ball is
accepted only if it is more than one ring away from every carved ball, its
boundary is a $2$-sphere, and it meets the boundary of the cube in no triangle (a
cavity: $\beta_2$ grows by one) or in one disk (a dent: no change); other
contacts, which could bore a tunnel, are rejected. Refilling is the exact
reverse of a carving. The number of carved balls follows the same walk as the
number of holes in W1, with target $4$. The Betti numbers are $(1,0,c,0)$, where
$c$ is the number of cavities.

\paragraph{W4: flag complexes of moving points.}
$n$ points are drawn uniformly in the unit square or the unit cube (``2D'' and
``3D'' in the tables and figures). A radius
$\varepsilon$ is calibrated on the initial points so that the
$\varepsilon$-neighbourhood graph has exactly $5n$ edges (mean degree $10$), and
$K$ is the flag complex of this graph truncated at dimension $D\in\{2,3\}$. An
update moves one uniformly random point by adding to each coordinate an
independent offset drawn uniformly from $[-\varepsilon/4,\varepsilon/4]$,
reflected at the border of the domain, and keeps its vertex identifier: it
deletes the closed star of the vertex top-down and inserts the vertex and its new
star bottom-up. The complexes are not manifolds and all Betti numbers vary; they
are not known in closed form and are checked against the oracles only.

\emph{Caveat: truncation.} In a flag complex truncated at dimension $D$, the
vertices of a clique of size $D+2$ span a $(D+1)$-simplex that is not in~$K$, so
its boundary is a $D$-cycle that is never filled. Hence $\beta_D$ is large: at
INIT, $\beta_D/n$ lies between \NumWFourBetaDRatioMin{} and
\NumWFourBetaDRatioMax{} in our W4 instances (\cref{tab:w4-betti}). This is genuine linear algebra on
a non-manifold complex, but $\beta_D$ is dominated by the truncation, not by the
shape of the point cloud. We use W4 as an algebraic stress test and draw no
conclusion about the topology of point clouds from it.

\paragraph{W5: non-locality witnesses.}
(a) The cycle graph $C_L$, whose updates toggle the edge $\{0,1\}$: odd-numbered
updates delete it and even-numbered ones insert it again. (b) The icosphere of
frequency $2^L$, whose updates toggle one triangle chosen from the seed in the
same way. By Proposition~\ref{main-prop:nonlocal}, every toggle costs every DRB method
$\Omega(L)$, respectively $\Omega(n_2)$, operations; once the states alternate,
a toggle costs $O(L+\nnz(R)+\nnz(V))$, respectively $O(n_2+\nnz(R)+\nnz(V))$.
\method{euler} needs $O(1)$ per toggle on (b) and does not apply to (a), which
is a graph and not a surface.

\paragraph{Sweeps.}
\Cref{tab:sweeps} lists the workloads, parameters and methods of each
experiment. Every timed configuration ran with \NumSeeds{} seeds.

\begin{table}[t]
  \centering
  \caption{Experiments. Parameters not listed are the generator defaults of
    \cref{sec:exp-workloads}. ``as E1'': \method{full-id}, \method{full-sorted}, \method{full-morse},
    \method{gudhi}, \method{drb-id}, \method{drb-id-rb}, \method{mmhm},
    \method{mmhm-norb}; ``DRB methods'': the last four. \method{gudhi} processes
    all updates on instances with at most \NumGudhiFullLimit{} simplices and the
    first $T'={}$\NumPrefix{} updates on larger ones.}
  \label{tab:sweeps}
  \small
  \begin{tabularx}{\linewidth}{@{}l >{\raggedright\arraybackslash}p{3.2cm} >{\raggedright\arraybackslash}X >{\raggedright\arraybackslash}p{4.3cm}@{}}
    \toprule
    & Workload & Sweep & Methods \\
    \midrule
    E0 & samples; all runs of E1--E8 & --- & all exact methods \\
    E1 & W2, $s=8$ & $N\in\{7,10,15,22,33\}$ & as E1 \\
    E2 & W2, $N=22$ & $s\in\{1,2,4,8,16,32,64\}$ & DRB methods, \method{full-id}, \method{full-morse} \\
    E3 & W3, $k=2$ & $N\in\{14,20,27\}$ & as E1 \\
    E4 & W4 & square, $D=2$: $n\in\{1000,4000,16000\}$; square and cube, $D=3$: $n\in\{1000,4000\}$ & as E1 \\
    E5 & W1 & sphere $L\in\{4,5,6,7\}$; torus $n=m\in\{64,128,256,500\}$ & as E1, and \method{euler} \\
    E6 & W5a, W5b & (a) $L=10^3,3\cdot10^3,10^4,\dots,10^6$; (b) levels $L=3,\dots,8$ & DRB methods, \method{full-id}, \method{gudhi}; \method{euler} on (b) \\
    E7 & W2 $N=22$, W3 $N=20$, W4 square $n=16000$, $D=2$ & $\alpha\in\{0.25,0.5,1,2,4,\infty\}$ \mbox{$\times$} $m\in\{\infty,256,64\}$ & \method{mmhm}, \method{drb-id-rb} \\
    E8 & W2 $N=22$, W3 $N=20$, W4 as in E7, W1 torus $256\times256$; for \method{drb-rand} one size per family (\cref{tab:rand}) & runs of E1, E3, E4, E5 & DRB methods; \method{drb-rand} \\
    \bottomrule
  \end{tabularx}
\end{table}

\subsection{Timing protocol and drift of the static baseline}\label{sec:exp-timing}


The input sequence is parsed completely before timing starts. For the engine
methods, the time of update $t$ is the wall-clock time
(\texttt{std::chrono::steady\_clock}) of processing its block of atomic edits:
the store updates, the work of the method for each edit, the recomputation of
the static baselines, and any rebuild, which is charged to the update that
triggers it. Verification and output are excluded. Building the structure for
$K_0$ is timed separately and reported as build time. For \method{gudhi}, the
time of an update covers the simplex-tree edits, \texttt{compute\_persistence}
and \texttt{betti\_numbers}, again without parsing.

All runs are single-threaded processes. Each configuration ran with
\NumSeeds{} seeds, which change the generated workload; W5a is determined by $L$
alone, so its three runs are repetitions of one instance. Every run processed
$T={}$\NumUpdates{} updates, except \method{gudhi} on instances with more than
\NumGudhiFullLimit{} simplices, which processed the first $T'={}$\NumPrefix{}
(hollow markers in the figures, $'$ in the tables). For each instance and method
we pool the per-update latencies of the three runs and report their median,
interquartile range and 95th and 99th percentiles (p95, p99), and the
\emph{amortised latency}: the total processing time of all updates, rebuilds
included, divided by their number (the initial build is not included; we
report it separately). A \emph{speedup} is a ratio of amortised
latencies over the same updates; for \method{gudhi} on a prefix, both methods
are taken over that prefix. Slopes are least-squares fits of $\log$ latency
against $\log$ size. All
\NumConfiguredRuns{} configured runs finished (\NumNotRun{} were not run); they
took \NumComputeHours{} hours of single-threaded processing in total.

\paragraph{Drift of the static baseline.}
Because the insertion order of \method{full-id} depends on the history, its cost
is not stationary. Its median latency over the last 200 updates exceeds that
over the first 200 by more than 25\,\% on \NumDriftFullIdInstances{} instances:
every W3 instance (up to a factor \NumWThreeNTwentySevenFullIdDriftMed), every
W4 instance with $n<16000$, and the $64\times64$ torus (\cref{tab:drift}).
\method{full-sorted}, whose order cannot drift, drifts by at most
\NumWOneSphereLFourFullSortedDriftMed{} on W1 and W2, by
\NumWFourSpaceNOneThousandDThreeFullSortedDriftMed{} to
\NumWFourPlaneNFourThousandDThreeFullSortedDriftMed{} on W4 (\method{full-id}:
\NumWFourPlaneNSixteenThousandDTwoFullIdDriftMed{} to
\NumWFourPlaneNOneThousandDTwoFullIdDriftMed), and by
\NumWThreeNFourteenFullSortedDriftMed{} to
\NumWThreeNTwentyFullSortedDriftMed{} on W3, where part of the drift therefore
comes from the changing complex rather than from the order. The DRB methods drift
as well; on W3 \method{mmhm} drifts less than \method{full-id}. Where the drift
matters we give speedups over matched windows of the first and the last 200
updates and against \method{full-sorted}.


\paragraph{Counters.}
Besides time, every run records per update the number of column additions, the
sum of the footprints $|\occ(\sigma)|$ over the deleted simplices, and
$\nnz(R)+\nnz(V)$; in an update with a rebuild, the additions of the rebuild and
of the edits are summed. These counts are proxies for the cost terms of
Proposition~\ref{main-prop:cost}: they count terms, not the sizes of the chains involved
(Section~\ref{main-sec:cost}), so they need not move with time. They do not depend on the
machine.


\section{Noise and an outlier run}\label{sec:s-noise}

\paragraph{Noise.}
Five identical runs of \method{mmhm} and \method{full-id} on W2 with $N=22$ gave
coefficients of variation (CV) of the amortised latency of
\NumNoiseCVMmhmAmort{} and \NumNoiseCVFullIdAmort\,\% and of the p99 of
\NumNoiseCVWTwoNTwentyTwoSEightMmhmPNinetyNine{} and
\NumNoiseCVWTwoNTwentyTwoSEightFullIdPNinetyNine\,\%. On the W1 torus
$256\times256$, W3 with $N=20$ and W5a with $L=10^5$ (five runs each of
\method{mmhm}, \method{drb-id} and \method{full-id}) the amortised CV was
\NumNoiseCVWOneTorusTwoHundredFiftySixDrbIdAmort{} to
\NumNoiseCVWThreeNTwentyMmhmAmort\,\% and the p99 CV
\NumNoiseCVWOneTorusTwoHundredFiftySixMmhmPNinetyNine{} to
\NumNoiseCVWFiveCycleLOneHundredThousandFullIdPNinetyNine\,\%; on
microsecond-scale updates (W5a, $L=1000$) the CV reaches
\NumNoiseCVAmortMax\,\% (amortised) and \NumNoiseCVPNinetyNineMax\,\% (p99). At
most two timed runs executed at once. Concurrent load changed the median within
the run-to-run spread at $N=22$, but at \NumWTwoNThirtyThreeSEightSize{} simplices
three busy processes raised it by $\NumNoiseFullIdLargeFourBusyPct\,\%$
(\method{full-id}) and $\NumNoiseMmhmLargeFourBusyPct\,\%$ (\method{mmhm});
untimed verification jobs overlapped \NumTimedRunsWithUntimed{} timed runs early
in the sweep (\path{experiments/ENVIRONMENT.md}). We call a difference between
two methods \emph{within noise} when their per-seed ranges overlap
(\path{experiments/results/summary_pairs.csv}). Each range has three points, and
on W1--W4 the seeds are different instances, so this is a weak criterion, not a
test.


The outlier run is described in \cref{sec:e1}; \cref{tab:noise} lists the
repeated runs.

\begin{table}[ht]
  \centering
  \caption{Repeated identical runs: coefficients of variation of the median,
    amortised and p99 latency.}
  \label{tab:noise}
  \small
  \generatedtable{tab_noise}
\end{table}

\section{Detailed validation}\label{sec:s-validation}


The reported Betti numbers are compared with up to three references: (1) the
engine's \method{full-id}, a static recomputation that shares the parser and the
complex store with the dynamic methods but none of their reduction code, and is
therefore \emph{engine-internal}; (2) the closed-form Betti numbers of W1, W2, W3
and W5; and (3) GUDHI, an independent library in another language, which replays
the same sequence. The two external references are (2) and (3). A pure-Python
implementation of the dynamic reduced basis and of the Morse order, written from
the specification and sharing no code with the engine, ran only on the
\NumReferenceSamples{} committed samples.

Every exact run records the Betti numbers it reports after every update.
\method{full-id} processes all $T$ updates on every instance, and every other
exact run is compared state by state with the \method{full-id} output of the same
instance. The runs are also compared with the closed-form Betti numbers where
these are known (W1, W2, W3, W5), and with the GUDHI oracle on all updates of
instances with at most \NumGudhiFullLimit{} simplices and on the first $T'$
updates of larger ones. In addition, every exact engine run compares itself with
an internal recomputation after $K_0$, after every $50$th update and after the
last one, outside the timed region. On the \NumSamples{} committed samples, all
exact engine configurations (including aggressive rebuild settings), the
reference implementation, GUDHI and the closed-form values are cross-checked on
every state. A mismatch of an exact method stops the whole pipeline.
Table~\ref{main-tab:coverage} gives the resulting coverage.

Below the level of outputs, the engine's randomised tests check (I1)--(I3) and
compare every exact method with a brute-force rank computation on random edit
sequences, and the reference implementation checks (I1)--(I3) after every
atomic edit, Lemma~\ref{main-lem:apparent} and Corollaries~\ref{main-cor:critical-work} and~\ref{main-cor:morse-reduction} on random
complexes and Proposition~\ref{main-prop:nonlocal} on the W5 families.


% ---------------------------------------------------------------------------
% Results. One subsection per question of experiments/PLAN.md. Every number is
% a \Num macro (generated/numbers.tex) or sits in a generated table.
% ---------------------------------------------------------------------------


\subsection{Results (E0)}\label{sec:e0}


Every exact method reported correct Betti numbers in every checked state, and
every state was checked against at least one reference.

The \NumExactRuns{} runs of exact methods reported \NumCoverageVectors{} Betti
vectors (\cref{tab:correctness}). Every one was compared with at least one
reference, \NumCoverageGeTwoPct\,\% with at least two,
\NumCoverageExternalPct\,\% with at least one external reference (closed form or
GUDHI) and \NumCoverageAllThreePct\,\% with all three (Table~\ref{main-tab:coverage}); no
comparison found a mismatch (\NumExactMismatches). The
\NumCoverageEngineOnly{} vectors checked only against the engine-internal
\method{full-id} are those of the \method{gudhi} runs on W4, which has no closed
form. In addition, the exact engine runs made \NumInternalChecks{} internal
verifications against a recomputation, and on the
\NumReferenceSamples{} committed samples the Python reference checked
\NumReferenceChecks{} vectors of \NumReferenceMethods{} methods
(\NumReferenceMismatches{} mismatches).

\begin{table}[t]
  \centering
  \caption{Correctness (E0): per exact method, the runs, the Betti vectors they
    reported, the states compared with the \method{full-id} output of the same
    instance, with the closed-form values and with GUDHI, and the mismatches
    (\method{drb-rand} ran on four instances only).}
  \label{tab:correctness}
  \small
  \generatedtable{tab_correctness_exact}
\end{table}

\section{Complete performance results}\label{sec:s-results}

\subsection{W2: dependence on the size of the complex (E1)}\label{sec:e1}


Only weakly for the dynamic reduced basis, and linearly for static
recomputation. On W2 with $s=8$, the median per-update latency of the four DRB
methods grows with $|K_0|$ with slope \NumSlopeEOneWTwoSEightMmhmMed{} to
\NumSlopeEOneWTwoSEightDrbIdMed{} over five sizes, against
\NumSlopeEOneWTwoSEightFullIdMed{} for \method{full-id},
\NumSlopeEOneWTwoSEightGudhiMed{} for \method{gudhi} and
\NumSlopeEOneWTwoSEightFullMorseMed{} for \method{full-morse}
(\cref{fig:full-scaling}(a)).

At the largest size, $|K_0|={}$\NumWTwoNThirtyThreeSEightSize, the four DRB
variants are \NumWTwoNThirtyThreeSEightDrbSpeedupFullIdMin{} to
\NumWTwoNThirtyThreeSEightDrbSpeedupFullIdMax{} times faster than
\method{full-id}, and their differences are within noise: no rebuild fired, and
the insertion order is as fast as the Morse order. On W2 the gain is that of the
dynamic reduced basis, not of the Morse order or of rebuilds. \method{mmhm} takes
\NumWTwoNThirtyThreeSEightMmhmAmort\,\textmu s per update amortised against
\NumWTwoNThirtyThreeSEightFullIdAmort\,\textmu s for \method{full-id}, and its
p99 is \NumWTwoNThirtyThreeSEightMmhmPNinetyNine{} against
\NumWTwoNThirtyThreeSEightFullIdPNinetyNine\,\textmu s
(\cref{tab:summary-largest,tab:full-p99}). Over the same first $T'$ updates,
\method{mmhm} is \NumWTwoNThirtyThreeSEightMmhmSpeedupGudhi{} times faster than
\method{gudhi}. The DRB pays more at the start of a run: over the first 200
updates \method{mmhm} is \NumWTwoNThirtyThreeSEightMmhmSpeedupFullIdFirstBlock{}
times faster than \method{full-id}, over the last 200
\NumWTwoNThirtyThreeSEightMmhmSpeedupFullIdLastBlock{} times, although its median
hardly changes (\NumWTwoNThirtyThreeSEightMmhmFirstBlockMed{} and
\NumWTwoNThirtyThreeSEightMmhmLastBlockMed\,\textmu s); a few early updates are
expensive. \method{full-sorted}, which does not drift on W2 and is within noise
of \method{full-id} at this size, gives the same picture: \method{mmhm} is
\NumWTwoNThirtyThreeSEightMmhmSpeedupFullSorted{} times faster over the whole
run, \NumWTwoNThirtyThreeSEightMmhmSpeedupFullSortedFirstBlock{} and
\NumWTwoNThirtyThreeSEightMmhmSpeedupFullSortedLastBlock{} times over the two
windows. Building the structure for $K_0$, including the insertion into the
store, took \NumWTwoNThirtyThreeSEightMmhmInitMs\,ms for \method{mmhm} and
\NumWTwoNThirtyThreeSEightFullIdInitMs\,ms for \method{full-id}, and
$\nnz(R)+\nnz(V)$, a proxy for the memory of the DRB, reached at most
\NumWTwoNThirtyThreeSEightMmhmNnzPerSimplex{} entries per simplex.

One run is an outlier: at $N=22$, one of the three \method{mmhm-norb} runs was
slow in its first 400 updates only (median \NumOutlierSeedOneBlockOneMed\,\textmu s,
against \NumOutlierSeedZeroBlockOneMed{} and \NumOutlierSeedTwoBlockOneMed\,\textmu s
for the other two). A rerun produced identical Betti numbers, additions, $\nnz$
and footprints and normal timings (amortised \NumOutlierRerunAmort{} against
\NumOutlierOrigAmort\,\textmu s); the tables keep the original run.

\paragraph{Static recomputation in the Morse order.}
\method{full-morse} is slower than \method{full-id} and \method{full-sorted} on
every instance on which they ran, than \method{full-id} by factors from \NumWOneSphereLFourFullMorseOverFullId{} (W1 sphere,
$L=4$) to \NumWTwoNThirtyThreeSEightFullMorseOverFullId{} (W2, $N=33$;
\cref{tab:full-amort}). There is little for the collapse to save: on W2 the
insertion order needs only \NumWTwoNThirtyThreeSEightFullIdAdds{} additions per
recomputation, and on the largest W4 instance \method{full-morse} performs more
additions than \method{full-id} (\NumWFourPlaneNSixteenThousandDTwoFullMorseAdds{}
against \NumWFourPlaneNSixteenThousandDTwoFullIdAdds). For one-shot static
recomputation the collapse does not pay for itself on our workloads.

\begin{figure}[t]
  \centering
  \resultfigure[width=0.86\linewidth]{F2_latency_vs_size}{5.5cm}
  \caption{Median per-update latency against $|K_0|$ at a fixed edit size
    (log--log; band: interquartile range; three runs pooled). (a)~W2 with $s=8$
    (E1); (b)~W3 (E3); (c)~W4 in the square with $D=2$ (E4); (d)~W1 punctured
    sphere (E5), without the bands of the DRB methods, which overlap. Hollow
    markers: \method{gudhi} on the prefix $T'$.}
  \label{fig:full-scaling}
\end{figure}

\begin{table}[tp]
  \centering
  \caption{Summary at the largest size of each family of E1, E3, E4 and E5:
    $|K_0|$, the number $T$ of timed updates per run, the median, p95 and
    amortised per-update latency in \textmu s (three runs pooled), and the
    speedup over \method{full-id} (ratio of amortised latencies over all $T$
    updates). \method{gudhi} entries with $T={}$\NumPrefix{} cover only the
    first $T'$ updates, and on W3 \method{full-id} drifts (\cref{tab:drift}); the text gives
    matched comparisons. Every instance is listed in
    \cref{tab:full-amort,tab:full-p99}.}
  \label{tab:summary-largest}
  \small
  \generatedtable{tab_summary}
\end{table}

\subsection{W2: dependence on the edit size (E2)}\label{sec:e2}


Linearly for the DRB methods. On W2 with $N=22$ ($|K_0|$ between
\NumWTwoNTwentyTwoSOneSize{} and \NumWTwoNTwentyTwoSSixtyFourSize), the median
per-update latency of the DRB methods grows with the number $s$ of stellar moves
with slope \NumSlopeETwoWTwoNTwentyTwoDrbIdRbMed{} to
\NumSlopeETwoWTwoNTwentyTwoMmhmMed, whereas \method{full-id} and
\method{full-morse} barely depend on $s$ (slopes
\NumSlopeETwoWTwoNTwentyTwoFullIdMed{} and
\NumSlopeETwoWTwoNTwentyTwoFullMorseMed; \cref{fig:edit-size}(a)). The amortised
cost per atomic edit of the DRB methods therefore stays in a narrow band, while
that of recomputation falls roughly as $1/s$ (\cref{fig:edit-size}(b)). There is
no crossover in the measured range: at $s=64$, that is $1024$ atomic edits per
update, the DRB methods are still \NumWTwoNTwentyTwoSSixtyFourMmhmSpeedupFullId{}
to \NumWTwoNTwentyTwoSSixtyFourDrbIdRbSpeedupFullId{} times faster than
\method{full-id}.

\begin{figure}[t]
  \centering
  \resultfigure[scale=1]{F3_latency_vs_edit_size}{4.5cm}
  \caption{Cost against the number $s$ of stellar moves per update (E2; W2,
    $N=22$, $16s$ atomic edits per update; log--log). (a)~Median per-update
    latency with interquartile range; (b)~amortised latency per atomic edit. The
    hollow \method{mmhm-norb} point at $s=8$ includes the outlier run of E1;
    $\times$ marks the same point with that run replaced by its rerun.}
  \label{fig:edit-size}
\end{figure}

\subsection{W3: cavities (E3)}\label{sec:e3}


\method{mmhm} has the lowest amortised latency at every size of W3, beyond
noise, but the DRB is not uniformly better than recomputation. We take
\method{full-sorted} as the reference here: it is faster than \method{full-id}
at every size and drifts less on the larger instances (\cref{sec:exp-timing}).
Against it, the DRB methods without rebuilds are slower at $N=14$ and have a
higher p99 at $N\le20$; the rebuilding ones are faster at every size.

At $N=14$ ($|K_0|={}$\NumWThreeNFourteenSize), \method{mmhm} takes
\NumWThreeNFourteenMmhmAmort\,\textmu s amortised, \method{full-sorted}
\NumWThreeNFourteenFullSortedAmort{} and \method{full-id}
\NumWThreeNFourteenFullIdAmort\,\textmu s; \method{drb-id}
(\NumWThreeNFourteenDrbIdAmort\,\textmu s) and \method{mmhm-norb}
(\NumWThreeNFourteenMmhmNorbAmort\,\textmu s) are slower than
\method{full-sorted}, and at $N=20$ \method{drb-id} still is (speedup
\NumWThreeNTwentyDrbIdSpeedupFullSorted). The p99 of \method{full-sorted},
\NumWThreeNFourteenFullSortedPNinetyNine\,\textmu s, is below that of every DRB
method at $N=14$ except \method{mmhm} (\NumWThreeNFourteenMmhmPNinetyNine\,\textmu s,
within noise). At $N=27$ ($|K_0|={}$\NumWThreeNTwentySevenSize) every DRB method
is faster than both static baselines on average. \method{mmhm} is
\NumWThreeNTwentySevenMmhmSpeedupFullSorted{} times faster than
\method{full-sorted} over the whole run
(\NumWThreeNTwentySevenMmhmSpeedupFullSortedFirstBlock{} and
\NumWThreeNTwentySevenMmhmSpeedupFullSortedLastBlock{} over the first and the
last 200 updates) and \NumWThreeNTwentySevenMmhmSpeedupFullId{} times faster
than \method{full-id}: \method{full-id} is
\NumWThreeNTwentySevenFullIdSpeedupFullSortedFirstBlock{} times as fast as
\method{full-sorted} over the first 200 updates and
\NumWThreeNTwentySevenFullIdSpeedupFullSortedLastBlock{} times as fast over the
last 200, as its order drifts (\cref{tab:drift}). Over the same first $T'$ updates \method{mmhm} is
\NumWThreeNTwentySevenMmhmSpeedupGudhi{} times faster than \method{gudhi}. At
$N=27$ the p99 of \method{mmhm} and \method{drb-id-rb} is below that of
\method{full-sorted}, that of \method{mmhm-norb} within noise of it, and that of
\method{drb-id} above it. At $N=20$ and $N=27$ the amortised latencies are
ordered \method{mmhm}, \method{drb-id-rb}, \method{mmhm-norb}, \method{drb-id},
all beyond noise (at $N=14$ the last two are within noise): rebuilding matters
more than the order, and the Morse order helps both with and without rebuilds.

The basis grows over a W3 run (\cref{fig:timeseries}(b); it also grows on W1,
E5). At $N=20$, without rebuilds, $\nnz(R)+\nnz(V)$ grows from
\NumWThreeNTwentyDrbIdNnzStart{} to \NumWThreeNTwentyDrbIdNnzEnd{}
(\method{drb-id}) and to \NumWThreeNTwentyMmhmNorbNnzEnd{} (\method{mmhm-norb}),
and the latency grows with it. The growth trigger rebuilds \method{mmhm}
\NumWThreeNTwentyMmhmRebuilds{} times and \method{drb-id-rb}
\NumWThreeNTwentyDrbIdRbRebuilds{} times per run (means over the runs), and the
runs end at \NumWThreeNTwentyMmhmNnzEnd{} and \NumWThreeNTwentyDrbIdRbNnzEnd.

\begin{figure}[t]
  \centering
  \resultfigure[width=0.82\linewidth]{F4_time_series}{6cm}
  \caption{Behaviour over time of the four DRB methods (seed~0). Top: per-update
    latency, mean of a sliding window of 50 updates, with the rebuilds of
    \method{mmhm} and \method{drb-id-rb} marked. Bottom: $\nnz(R)+\nnz(V)$ after
    each update. (a)~W2 with $N=22$, $s=8$ (E1); (b)~W3 with $N=20$ (E3);
    (c)~W1 torus $256\times256$ (E5), with the updates after which no hole is
    open marked.}
  \label{fig:timeseries}
\end{figure}

\subsection{W4: flag complexes (E4)}\label{sec:e4}


The DRB methods are faster than both static baselines on every W4 instance, their
latency hardly depends on $|K_0|$, and on the larger instances the Morse order
helps. In the square with $D=2$ (\cref{fig:full-scaling}(c)), their median
latency grows with slope \NumSlopeEFourWFourPlaneDTwoMmhmNorbMed{} to
\NumSlopeEFourWFourPlaneDTwoDrbIdRbMed, against
\NumSlopeEFourWFourPlaneDTwoFullIdMed{} for \method{full-id} (three sizes, so the
slopes are indicative only). Against \method{full-sorted}, our reference here
because it drifts less than \method{full-id} on every W4 instance
(\cref{tab:drift}), their speedup is
\NumWFourSpaceNOneThousandDThreeDrbIdRbSpeedupFullSorted{} to
\NumWFourPlaneNOneThousandDTwoMmhmNorbSpeedupFullSorted{} at $n=1000$ and
\NumWFourPlaneNFourThousandDThreeDrbIdSpeedupFullSorted{} to
\NumWFourPlaneNSixteenThousandDTwoMmhmNorbSpeedupFullSorted{} at $n\ge4000$;
against \method{full-id} it is
\NumWFourSpaceNOneThousandDThreeDrbIdRbSpeedupFullId{} to
\NumWFourPlaneNOneThousandDTwoMmhmNorbSpeedupFullId{} and
\NumWFourPlaneNFourThousandDThreeDrbIdSpeedupFullId{} to
\NumWFourPlaneNSixteenThousandDTwoMmhmNorbSpeedupFullId. At $n=16000$
\method{full-sorted} is slower than \method{full-id}
(\NumWFourPlaneNSixteenThousandDTwoFullSortedAmort{} against
\NumWFourPlaneNSixteenThousandDTwoFullIdAmort\,\textmu s amortised), and
\method{mmhm} is \NumWFourPlaneNSixteenThousandDTwoMmhmSpeedupFullSorted{} times
faster than it over the whole run and
\NumWFourPlaneNSixteenThousandDTwoMmhmSpeedupFullSortedFirstBlock{} and
\NumWFourPlaneNSixteenThousandDTwoMmhmSpeedupFullSortedLastBlock{} times over the
first and the last 200 updates (against \method{full-id}:
\NumWFourPlaneNSixteenThousandDTwoMmhmSpeedupFullId,
\NumWFourPlaneNSixteenThousandDTwoMmhmSpeedupFullIdFirstBlock{} and
\NumWFourPlaneNSixteenThousandDTwoMmhmSpeedupFullIdLastBlock). The p99 of the
DRB methods is below that of both baselines on every W4 instance except the two
with $n=1000$ and $D=3$, where the difference is within noise
(\cref{tab:full-p99}).

\method{mmhm-norb} is faster than \method{drb-id} beyond noise on the square
instances with $n\ge4000$; on the others the difference is within noise. At
$n=16000$ the amortised latencies are
\NumWFourPlaneNSixteenThousandDTwoMmhmNorbAmort{} and
\NumWFourPlaneNSixteenThousandDTwoDrbIdAmort\,\textmu s, the p99 latencies
\NumWFourPlaneNSixteenThousandDTwoMmhmNorbPNinetyNine{} and
\NumWFourPlaneNSixteenThousandDTwoDrbIdPNinetyNine\,\textmu s, and the p99 of the
column additions per update
\NumEEightWFourPlaneNSixteenThousandDTwoMmhmNorbAddsPNinetyNine{} and
\NumEEightWFourPlaneNSixteenThousandDTwoDrbIdAddsPNinetyNine{} (E8). Under the
default policy, rebuilds fired in only one W4 run of each rebuilding method (in
the cube, $n=1000$, $D=3$). Without rebuilds, \method{mmhm} does the same
additions as \method{mmhm-norb} (\cref{tab:e8}) but is slower at $n=16000$
(\NumWFourPlaneNSixteenThousandDTwoMmhmAmort{} against
\NumWFourPlaneNSixteenThousandDTwoMmhmNorbAmort\,\textmu s, beyond noise). We have
not profiled the difference. One difference between the two is the index of
cells by vertex tuple that \method{mmhm} maintains for a possible Morse
rebuild; in E7, \method{mmhm} with both triggers off, which never rebuilds but
keeps the index, has a median of
\NumESevenWFourPlaneNSixteenThousandDTwoMmhmAOffPOffMed\,\textmu s against
\NumWFourPlaneNSixteenThousandDTwoMmhmNorbMed\,\textmu s for \method{mmhm-norb}.


\section{Surfaces (W1, E5)}\label{sec:s-surfaces}\label{sec:e5}


On average the DRB is at best about twice as fast as recomputation, sometimes
slower, and it has the worse tail; the combinatorial counters are far cheaper
than both. On no W1 instance is a DRB method more than
\NumWOneSphereLFiveMmhmSpeedupFullId{} times faster than \method{full-id}, and on
several smaller instances \method{mmhm-norb} is within noise of it. Against
\method{full-sorted}, \method{drb-id} is slower on every W1 instance (within noise
on the $256\times256$ torus), \method{mmhm-norb} is slower on the four smallest,
and no DRB method is more than \NumWOneSphereLSevenMmhmNorbSpeedupFullSorted{}
times faster. On the
largest sphere and torus the insertion-order variants are slower than
\method{full-id} and the Morse-ordered ones faster, all beyond noise. Every DRB
variant has a higher p99 than \method{full-id} on every W1 instance.

On the sphere with $L=7$ ($|K_0|={}$\NumWOneSphereLSevenSize), \method{full-id}
takes \NumWOneSphereLSevenFullIdAmort\,\textmu s amortised, \method{drb-id}
\NumWOneSphereLSevenDrbIdAmort{} and \method{drb-id-rb}
\NumWOneSphereLSevenDrbIdRbAmort\,\textmu s; \method{mmhm-norb} and
\method{mmhm} are \NumWOneSphereLSevenMmhmNorbSpeedupFullId{} and
\NumWOneSphereLSevenMmhmSpeedupFullId{} times faster than \method{full-id}. On
the torus with $n=m=500$ ($|K_0|={}$\NumWOneTorusFiveHundredSize) the speedups
are \NumWOneTorusFiveHundredDrbIdSpeedupFullId{} (\method{drb-id}),
\NumWOneTorusFiveHundredDrbIdRbSpeedupFullId{} (\method{drb-id-rb}),
\NumWOneTorusFiveHundredMmhmSpeedupFullId{} (\method{mmhm}) and
\NumWOneTorusFiveHundredMmhmNorbSpeedupFullId{} (\method{mmhm-norb}). The median
tells a different story from the average: at $L=7$ the median of \method{mmhm}
is \NumWOneSphereLSevenMmhmSpeedupFullIdMed{} times below that of
\method{full-id} (\NumWOneSphereLSevenMmhmMed{} against
\NumWOneSphereLSevenFullIdMed\,\textmu s), while its p99 is
\NumWOneSphereLSevenMmhmPNinetyNine{} against
\NumWOneSphereLSevenFullIdPNinetyNine\,\textmu s; the average is set by a heavy
tail (\cref{tab:summary-largest,tab:full-p99}). \method{euler} needs
\NumWOneSphereLFourEulerAmort{} to \NumWOneTorusFiveHundredEulerAmort\,\textmu s
per update, including the edits of the complex store,
\NumWOneSphereLSevenEulerSpeedupFullId{} times less than \method{full-id} on the
largest sphere.

The median of the DRB methods grows with $|K_0|$ almost like that of
\method{full-id} (slopes \NumSlopeEFiveWOneSphereMmhmNorbMed{} to
\NumSlopeEFiveWOneSphereDrbIdRbMed{} on the sphere, against
\NumSlopeEFiveWOneSphereFullIdMed; \cref{fig:full-scaling}(d)); the chain sizes
behind this were not logged. The closed surface,
on which Proposition~\ref{main-prop:nonlocal} forces $\Omega(n_2)$ work, reappears after only
\NumWOneSphereLSevenClosedStates{} of the \NumUpdates{} updates of a run on the
largest sphere (mean over the runs), too rarely to set the median, and on that
sphere the median of \method{mmhm-norb} is nearly the same over the first and the
last 200 updates (\NumWOneSphereLSevenMmhmNorbFirstBlockMed{} and
\NumWOneSphereLSevenMmhmNorbLastBlockMed\,\textmu s). The cause of the growth is open. Between rebuilds the basis also becomes denser
(\cref{fig:timeseries}(c)): on the largest torus $\nnz(R)+\nnz(V)$ grows over a
run by a factor \NumWOneTorusFiveHundredDrbIdNnzGrowth{} for \method{drb-id} and
\NumWOneTorusFiveHundredMmhmNorbNnzGrowth{} for \method{mmhm-norb}. This is
Remark~\ref{main-rem:surfaces} in practice: on connected surfaces with known connectivity,
the counters are correct and much cheaper, and the algebra is not needed.


\section{Non-locality witnesses (W5, E6)}\label{sec:s-nonlocal}\label{sec:e6}


Where Proposition~\ref{main-prop:nonlocal} says it does. On both witness families, the cost of
every DRB method per toggle grows linearly with the size of the toggled cycle.
The DRB methods stay faster than \method{full-id} on average, but their p99 is
higher on the cycle graph, and for the Morse-ordered variants also on the
largest sphere; only \method{euler} (on the sphere) is nearly independent of the
size.

The amortised latency per toggle grows with slope
\NumSlopeESixWFiveCycleDrbIdAmort{} to \NumSlopeESixWFiveCycleDrbIdRbAmort{}
on W5a ($L$ from $10^3$ to $10^6$) and \NumSlopeESixWFiveSphereDrbIdRbAmort{} to
\NumSlopeESixWFiveSphereMmhmAmort{} on W5b (levels 3 to 8), against
\NumSlopeESixWFiveCycleFullIdAmort{} and \NumSlopeESixWFiveSphereFullIdAmort{}
for \method{full-id} (Figure~\ref{main-fig:nonlocality}). At $L=10^6$, \method{drb-id}
takes \NumWFiveCycleLOneMillionDrbIdAmort\,\textmu s, \method{mmhm}
\NumWFiveCycleLOneMillionMmhmAmort\,\textmu s and \method{full-id}
\NumWFiveCycleLOneMillionFullIdAmort\,\textmu s; on the largest sphere
($|K_0|={}$\NumWFiveSphereLEightSize) the DRB methods take
\NumWFiveSphereLEightDrbIdAmort{} to \NumWFiveSphereLEightMmhmNorbAmort\,\textmu s
and \method{full-id} \NumWFiveSphereLEightFullIdAmort\,\textmu s. On W5a the p99
of every DRB method exceeds that of \method{full-id} at every $L$, beyond noise
except for \method{mmhm} at $L=1000$ (\cref{tab:full-p99}); at $L=10^6$ it is
\NumWFiveCycleLOneMillionMmhmPNinetyNine\,\textmu s for \method{mmhm} against
\NumWFiveCycleLOneMillionFullIdPNinetyNine\,\textmu s. On W5b at level~8 the p99
of \method{mmhm} and \method{mmhm-norb} is \NumWFiveSphereLEightMmhmPNinetyNineOverFullId{}
and \NumWFiveSphereLEightMmhmNorbPNinetyNineOverFullId{} times that of
\method{full-id}, that of the insertion-order variants
\NumWFiveSphereLEightDrbIdPNinetyNineOverFullId{} and
\NumWFiveSphereLEightDrbIdRbPNinetyNineOverFullId{} times. The Morse order does
not help here: at $L=10^6$, \method{drb-id} is faster than \method{mmhm} beyond
noise, and the reference implementation likewise counts more elementary
operations per toggle under the Morse order than under insertion order on these
families (\texttt{reference/README.md}); we have not identified why.
\method{euler} takes \NumWFiveSphereLThreeEulerAmort{} to
\NumWFiveSphereLEightEulerAmort\,\textmu s per toggle on W5b (slope
\NumSlopeESixWFiveSphereEulerAmort).


\section{Rebuild policy (E7)}\label{sec:s-rebuild}\label{sec:e7}


It depends on the workload, and no setting is best everywhere. On W3, where the
basis grows, rebuilding by the growth trigger pays off and the Morse order at a
rebuild helps. On W2 and W4 the trigger hardly ever fires and periodic rebuilds
only add cost (\cref{fig:rebuild}). On the surfaces of W1, where
the basis also grows, rebuilds help on the smaller instances and not on the
largest (\cref{tab:rebuild-effect}).

\paragraph{W2 and W4.} With the growth trigger alone, no rebuild fired on W2 at
any $\alpha$, and on W4 one \method{mmhm} run of three rebuilt once, at
$\alpha=0.25$ (\NumESevenWFourPlaneNSixteenThousandDTwoMmhmAQuarterPOffRebuilds{}
rebuilds per run). Periodic rebuilds raised the amortised latency of
\method{mmhm} on W2 from \NumESevenWTwoNTwentyTwoSEightMmhmAOnePOffAmort\,\textmu s
($\alpha=1$, no period) to
\NumESevenWTwoNTwentyTwoSEightMmhmAOnePTwoHundredFiftySixAmort\,\textmu s
($m=256$) and \NumESevenWTwoNTwentyTwoSEightMmhmAOnePSixtyFourAmort\,\textmu s
($m=64$). A Morse rebuild, which includes a collapse, costs more than an
insertion-order one: \method{drb-id-rb} took
\NumESevenWTwoNTwentyTwoSEightDrbIdRbAOnePSixtyFourAmort\,\textmu s at $m=64$.
\method{mmhm} is faster than \method{drb-id-rb} at only
\NumESevenWTwoNTwentyTwoSEightMmhmBelowDrbIdRb{} of the
\NumESevenWTwoNTwentyTwoSEightSettings{} settings on W2 and at
\NumESevenWFourPlaneNSixteenThousandDTwoMmhmBelowDrbIdRb{} on W4.

\paragraph{W3.} At $\alpha=1$ without a period, \method{mmhm} took
\NumESevenWThreeNTwentyMmhmAOnePOffAmort\,\textmu s amortised against
\NumESevenWThreeNTwentyMmhmAOffPOffAmort\,\textmu s without rebuilds, and
\method{drb-id-rb} \NumESevenWThreeNTwentyDrbIdRbAOnePOffAmort{} against
\NumESevenWThreeNTwentyDrbIdRbAOffPOffAmort\,\textmu s; the p99 of
\method{mmhm} fell from \NumESevenWThreeNTwentyMmhmAOffPOffPNinetyNine{} to
\NumESevenWThreeNTwentyMmhmAOnePOffPNinetyNine\,\textmu s. Too small a threshold
hurts again: at $\alpha=0.25$ \method{mmhm} rebuilt
\NumESevenWThreeNTwentyMmhmAQuarterPOffRebuilds{} times per run and its p99 was
\NumESevenWThreeNTwentyMmhmAQuarterPOffPNinetyNine\,\textmu s. On W3
\method{mmhm} is faster than \method{drb-id-rb} at
\NumESevenWThreeNTwentyMmhmBelowDrbIdRb{} of the
\NumESevenWThreeNTwentySettings{} settings. Under the default policy, the updates
after the first rebuild need fewer additions than the same updates of the
never-rebuilding twin (ratio \NumRebuildEffectWThreeNTwentyMmhmAfterFirst{} for
\method{mmhm}, \NumRebuildEffectWThreeNTwentyDrbIdRbAfterFirst{} for
\method{drb-id-rb} at $N=20$) and much less time
(\NumRebuildEffectWThreeNTwentyMmhmAfterFirstTime{} and
\NumRebuildEffectWThreeNTwentyDrbIdRbAfterFirstTime), on a much sparser basis
(\cref{fig:timeseries}(b)).

\paragraph{W1.} Under the default policy rebuilds fire on every W1 instance
(\cref{tab:rebuild-effect}). Against its never-rebuilding twin, \method{mmhm} is
faster beyond noise on the sphere with $L\le5$ and the tori with $n\le128$
(for instance \NumWOneSphereLFiveMmhmAmort{} against
\NumWOneSphereLFiveMmhmNorbAmort\,\textmu s at $L=5$), within noise on the
sphere with $L=7$ and the tori with $n\ge256$, and slower beyond noise on the
sphere with $L=6$. \method{drb-id-rb} is faster than \method{drb-id} on the
sphere with $L\le5$ and the tori with $n\le128$, within noise on the spheres
with $L\ge6$, and slower beyond noise on the tori with $n\ge256$. In the 100
updates after an \method{mmhm} rebuild, \method{mmhm} performs more additions
than \method{mmhm-norb} on the same updates on every W1 instance except the
sphere with $L=5$, by factors up to \NumRebuildEffectWOneTorusFiveHundredMmhmWindow.
Over all updates after the first rebuild, the ratio of additions is below one on
the large spheres (\NumRebuildEffectWOneSphereLSixMmhmAfterFirst{} and
\NumRebuildEffectWOneSphereLSevenMmhmAfterFirst) and above one on the largest
torus (\NumRebuildEffectWOneTorusFiveHundredMmhmAfterFirst), while the ratio of
time is \NumRebuildEffectWOneSphereLSixMmhmAfterFirstTime{} and
\NumRebuildEffectWOneSphereLSevenMmhmAfterFirstTime{} on those spheres. The
additions and the time disagree here, as the counters are only proxies
(Section~\ref{main-sec:cost}); we have not identified the mechanism.

\begin{figure}[t]
  \centering
  \resultfigure[width=0.82\linewidth]{F7_rebuild_sensitivity}{6cm}
  \caption{Rebuild policy (E7): amortised (top) and p99 (bottom) per-update
    latency of \method{mmhm} and \method{drb-id-rb} against the growth threshold
    $\alpha$, one line per period $m$; three runs pooled. (a)~W2 with $N=22$,
    $s=8$; (b)~W3 with $N=20$; (c)~W4 in the square with $n=16000$, $D=2$. With
    $\alpha=m=\infty$ (``off'') the methods never rebuild. E7 has no W1
    instance; W1 rebuild data are in \cref{tab:rebuild-effect}.}
  \label{fig:rebuild}
\end{figure}

\begin{table}[ht]
  \centering
  \caption{The full E7 grid: amortised and p99 latency and rebuilds per run of
    \method{mmhm} and \method{drb-id-rb} for every growth threshold $\alpha$ and
    period $m$.}
  \label{tab:e7}
  \scriptsize
  \setlength{\tabcolsep}{3pt}
  \generatedtable{tab_e7}
\end{table}

\section{What the Morse order buys (E8)}\label{sec:s-morse}\label{sec:e8}


At build time nothing against the insertion orders of our generators, although
both structured orders are far better than a random one. Between rebuilds it
depends on the
workload: the Morse order makes the DRB faster on surfaces, on cavities and on
the larger flag complexes, makes no difference on W2, and makes it slower on the
non-locality witnesses.

\paragraph{At build time.} Lemma~\ref{main-lem:apparent} and
Corollaries~\ref{main-cor:critical-work} and~\ref{main-cor:morse-reduction} describe the build, so we compare
the builds of \method{drb-id} and \method{mmhm-norb} (\cref{tab:build}). The
insertion orders written by our generators already need almost no additions on
W2 (\NumWTwoNThirtyThreeSEightDrbIdBuildAdds{} at $N=33$, against
\NumWTwoNThirtyThreeSEightMmhmNorbBuildAdds{} in the Morse order) and none on
W3; on W1 and W5 both orders need the same number within one percent (for
instance \NumWOneSphereLSevenDrbIdBuildAdds{} and
\NumWOneSphereLSevenMmhmNorbBuildAdds{} on the sphere with $L=7$). On W4 the
Morse order needs more additions, by factors
\NumWFourPlaneNOneThousandDThreeBuildAddsMorseOverId{} to
\NumWFourSpaceNFourThousandDThreeBuildAddsMorseOverId, and yields a slightly
smaller basis ($\nnz(R)+\nnz(V)$ ratio
\NumWFourSpaceNOneThousandDThreeNnzInitMorseOverId{} to
\NumWFourPlaneNOneThousandDTwoNnzInitMorseOverId). The lemma and its corollaries
thus say where the work of a build is done, not that there is less of it than
in a good insertion order.

\paragraph{Against a random order.} The insertion orders of our generators are
themselves structured. Built in a uniformly random order (\method{drb-rand}, one
instance per family, \cref{tab:rand}), the basis has
\NumWFourPlaneNFourThousandDTwoNnzInitRandOverId{} to
\NumWThreeNFourteenNnzInitRandOverId{} times as many entries after INIT as in
either structured order, and the build needs more additions on every instance
(on W2 \NumWTwoNFifteenSEightDrbRandBuildAdds, against
\NumWTwoNFifteenSEightDrbIdBuildAdds{} in insertion and
\NumWTwoNFifteenSEightMmhmNorbBuildAdds{} in Morse order). On W1, W2 and W3 the
updates are slower beyond noise, by factors
\NumWThreeNFourteenDrbRandOverDrbIdAmort{} to
\NumWTwoNFifteenSEightDrbRandOverDrbIdAmort{} against either structured order
(on W2 \NumWTwoNFifteenSEightDrbRandAmort{} against
\NumWTwoNFifteenSEightDrbIdAmort{} and
\NumWTwoNFifteenSEightMmhmNorbAmort\,\textmu s); there the choice between the two
structured orders matters less than using one. On W4 \method{drb-rand} is
within noise of \method{drb-id} and
\NumWFourPlaneNFourThousandDTwoDrbRandOverMmhmNorbAmort{} times slower than
\method{mmhm-norb}.

\paragraph{Between rebuilds.} Comparing the two methods that never rebuild,
\method{mmhm-norb} against \method{drb-id}, the Morse order is faster beyond
noise on the W1 spheres with $L\ge5$ and on all W1 tori (on the sphere
with $L=7$, \NumWOneSphereLSevenMmhmNorbAmort{} against
\NumWOneSphereLSevenDrbIdAmort\,\textmu s), on W3 with $N\ge20$ and on W4 in the
square with $n\ge4000$. It is within noise on W2 (except on the instance with
the outlier run of E1) and on the smaller W3 and W4 instances, and slower beyond noise on W5a for every $L\ge3000$ (at $L=10^6$,
\NumWFiveCycleLOneMillionMmhmNorbAmort{} against
\NumWFiveCycleLOneMillionDrbIdAmort\,\textmu s) and on W5b from level~6 on.
\Cref{fig:footprint,tab:e8} show the counters on the non-rebuild updates of one
instance each of W2, W3, W4 and W1; the footprint of an update is the sum of
$|\occ(\sigma)|$ over its deleted simplices. Without rebuilds the Morse order
gives lighter tails of the column additions on all four: on W4 ($n=16000$) the
p99 falls from \NumEEightWFourPlaneNSixteenThousandDTwoDrbIdAddsPNinetyNine{} to
\NumEEightWFourPlaneNSixteenThousandDTwoMmhmNorbAddsPNinetyNine{} and the median
footprint from \NumEEightWFourPlaneNSixteenThousandDTwoDrbIdFootprintMed{} to
\NumEEightWFourPlaneNSixteenThousandDTwoMmhmNorbFootprintMed. Not every quantile
improves: on W2 the median additions rise from
\NumEEightWTwoNTwentyTwoSEightDrbIdAddsMed{} to
\NumEEightWTwoNTwentyTwoSEightMmhmNorbAddsMed{} and the p99 footprint from
\NumEEightWTwoNTwentyTwoSEightDrbIdFootprintPNinetyNine{} to
\NumEEightWTwoNTwentyTwoSEightMmhmNorbFootprintPNinetyNine, and with rebuilds the
insertion-order rebuild of \method{drb-id-rb} leaves lighter footprint tails than
the Morse rebuild (p99 \NumEEightWThreeNTwentyDrbIdRbFootprintPNinetyNine{}
against \NumEEightWThreeNTwentyMmhmFootprintPNinetyNine{} on W3), although
\method{mmhm} has the lower amortised latency there (E7). Between rebuilds the
Morse order is a heuristic (Remark~\ref{main-rem:morse-scope}), and these measurements are
its whole justification.

\begin{figure}[t]
  \centering
  \resultfigure[width=0.82\linewidth]{F5_footprint_ccdf}{6cm}
  \caption{Counters between rebuilds (E8): complementary cumulative
    distributions of the footprint per update (top; the sum of
    $|\occ(\sigma)|$ over its deletions) and of the column additions per update
    (bottom), over the updates without a rebuild (three runs pooled; log--log).
    (a)~W2 with $N=22$, $s=8$; (b)~W3 with $N=20$; (c)~W4 in the square with
    $n=16000$, $D=2$. On W2 and W4 no rebuild fired, so \method{mmhm} coincides
    with \method{mmhm-norb} and \method{drb-id-rb} with \method{drb-id}. The W1
    torus of \cref{tab:e8} is not shown.}
  \label{fig:footprint}
\end{figure}

\section{Full result tables}\label{sec:s-tables}

Latencies are per update in \textmu s, the three runs pooled; $'$ marks
\method{gudhi} on the prefix $T'={}$\NumPrefix{}, and an empty cell a method not
run in that experiment.

\appendixtable{tab:build}{The build of $K_0$ in insertion order (\method{drb-id})
  and in the Morse order (\method{mmhm-norb}): $\nnz(R)+\nnz(V)$ after INIT and
  column additions of the build, with the ratios Morse/insertion (E8).}{tab_build}

\appendixtable{tab:drift}{Drift (\cref{sec:exp-timing}): median latency of
  \method{full-id} over the first and the last 200 updates and their ratio; the
  same ratio for \method{full-sorted}, \method{full-morse}, \method{drb-id} and
  \method{mmhm}; and the speedup of \method{mmhm} over \method{full-id} and over
  \method{full-sorted} on each of the two windows (ratio of their amortised
  latencies over the same updates).}{tab_drift}

\appendixtable{tab:rebuild-effect}{Does a rebuild help? Instances on which the
  default policy ($\alpha=1$, no period) rebuilt: rebuilds per run (rb:
  \method{drb-id-rb}, m: \method{mmhm}); ratios of amortised latency to the
  never-rebuilding twin (rb/id: \method{drb-id-rb} over \method{drb-id}, m/norb:
  \method{mmhm} over \method{mmhm-norb}); the same ratio of column additions over
  the 100 updates after each rebuild; and the ratios of additions and of time
  over all updates after the first rebuild. Above~1: the rebuilding method was
  costlier.}{tab_rebuild_effect}

\appendixtable{tab:full-med}{Median per-update latency (\textmu s) of every instance of E1--E6.}{tab_summary_full_med}

\appendixtable{tab:full-amort}{Amortised per-update latency (\textmu s) of every
  instance of E1--E6. The \method{mmhm-norb} entry for W2 with $N=22$, $s=8$ (in
  E1 and E2) includes the outlier run of \cref{sec:e1}. \method{gudhi} entries
  marked $'$ cover only the first $T'$ updates and are not comparable with the
  others.}{tab_summary_full}

\appendixtable{tab:e8}{Counters between rebuilds (E8): median and p99 of the
  column additions and of the footprint ($\sum|\occ(\sigma)|$ over the deletions)
  per update, over the updates without a rebuild, and $\nnz(R)+\nnz(V)$ after
  INIT and at the end (means over the runs).}{tab_e8}

\appendixtable{tab:full-p99}{99th-percentile per-update latency (\textmu s) of
  every instance of E1--E6. The \method{mmhm-norb} entry for W2 with $N=22$,
  $s=8$ includes the outlier run of \cref{sec:e1}.}{tab_summary_full_p99}

\appendixtable{tab:w4-betti}{W4: ranges of the Betti numbers over all states of
  the three runs, and $\beta_D/n$ at INIT. $\beta_D$ is dominated by unfilled
  boundaries of $(D+2)$-cliques (\cref{sec:exp-workloads}).}{tab_w4_betti}

\appendixtable{tab:rand}{Build order (E8): one instance per family built in
  insertion order (\method{drb-id}), in the Morse order (\method{mmhm-norb}) and
  in a random order (\method{drb-rand}): $\nnz(R)+\nnz(V)$ after INIT, column
  additions of the build, and amortised latency (three runs pooled).}{tab_rand}
